\selectlanguage{russian}
\section{Введение}
\label{sec:Introduction} \index{Introduction}

\subsection{Мотивация задачи}

Во многих прикладных задачах данные поступают группами, внутри которых наблюдаются две
совокупности совместно порождённых величин, но соответствие между их элементами утрачено.
Это происходит, например, при раздельном измерении двух модальностей, обезличивании записей
или независимом переупорядочивании результатов. В каждой группе известны только два набора
значений, тогда как неизвестно, какие элементы образуют пары.

\paragraph{Регистрация облаков точек.}
В компьютерном зрении и медицинской визуализации требуется совмещать облака точек,
описывающие один объект в разных системах координат. Отдельные точки двух облаков могут
соответствовать одним и тем же частям объекта, однако это соответствие непосредственно не
наблюдается. По нескольким таким группам требуется восстановить закономерность, связывающую
координаты в двух системах.

\paragraph{Кросс-модальный анализ одноклеточных данных.}
Различные молекулярные характеристики клетки нередко измеряются на разных клетках из одной
популяции. Отдельная выборка клеток образует группу, для которой доступны наборы измерений
двух модальностей, но неизвестно, какие значения относятся к одной клетке. Восстановленная
зависимость позволяет предсказывать одну молекулярную характеристику по другой.

Эти примеры различаются предметной областью, но приводят к общей статистической задаче:
по повторяющимся группам непарных наблюдений восстановить как скрытые пары, так и общую
зависимость между двумя величинами. При этом в обучающей выборке нет наблюдений с
достоверно известным попарным соответствием, по которым эту зависимость можно было бы
оценить непосредственно.

\subsection{Проблема существующих подходов}

Вероятностная постановка восстановления условного закона из групп непарных точек предложена
Байером и др.~\cite{beier2024transfer}. Неизвестное внутригрупповое соответствие в ней
трактуется как скрытая перестановка, а вычислительная аппроксимация строится с помощью
энтропийных транспортных ядер. Этот подход служит основной точкой сравнения в настоящей
работе.

Точное маргинальное правдоподобие модели со скрытой перестановкой содержит сумму по $M!$
допустимым сопоставлениям для группы размера $M$. Поэтому прямое вероятностное оценивание
становится вычислительно недоступным уже при умеренных $M$. Кроме того, транспортное
сопоставление для новой группы вычисляется заново и не использует обученный механизм
амортизированного вывода.

Дифференцируемые релаксации Gumbel--Sinkhorn~\cite{mena2018sinkhorn}, методы жёсткого
назначения~\cite{rezatofighi2018perm,carion2020detr} и релаксации сортировки
~\cite{prillo2020softsort,grover2019neuralsort} решают близкие комбинаторные задачи. Однако
они не дают непосредственно вероятностной модели условной плотности, обучаемой по группам
непарных наблюдений. Требуется связать оценивание условного распределения с вычислимым
вариационным выводом и ограничением, приближающим скрытую переменную к матрице перестановки.

\subsection{Цель и основные результаты}

Цель работы состоит в том, чтобы сделать вероятностное восстановление зависимости по
группам непарных данных вычислимым при росте размера групп и обеспечить обработку новых
групп без повторного решения отдельной задачи оптимизации. Для достижения этой цели в
работе получены следующие результаты.
\begin{enumerate}
  \item Сформулирована вариационная задача восстановления условного распределения при
  скрытом внутригрупповом сопоставлении и выведена соответствующая функция потерь на основе
  отрицательной нижней оценки обоснованности.
  \item Предложено использовать Gumbel--Sinkhorn-преобразование для дифференцируемого
  ограничения латентной матрицы назначений множеством мягких перестановок.
  \item Разработаны два подхода с разными параметризациями: амортизированный нейросетевой
  CatVAE, обучаемый стохастическим градиентным методом, и модель глобальной смеси сдвигов,
  обучаемая модифицированным EM с Sinkhorn-аппроксимацией E-шага и аналитическим M-шагом.
  \item Проведено сравнение с энтропийным транспортным оператором
  Байера и др.~\cite{beier2024transfer} по $L_2$-ошибке восстановления плотности при разных
  уровнях шума, размерах групп и объёмах данных.
\end{enumerate}
